\documentclass[../main]{subfiles}
\begin{document}
\ifSubfilesClassLoaded{\setcounter{page}{14}}{}
\section{Основа комунізму}
\label{sec:03_osnovanie}

Основою для капіталізму є товарне виробництво. Усе багатство капіталістичного суспільства виступає як сукупність товарів. Оскільки продукти праці за капіталізму неминуче набувають товарної форми, товар є «клітиною капіталізму».

Основою для комунізму як суспільного устрою, згідно з матеріалістичним розумінням історії, також може бути лише відповідний спосіб виробництва.

Усе багатство комуністичного суспільства – це універсальність змісту та форм людської діяльності. Марксизм взагалі бачить сутність людини в діяльності, спрямованій на зміну природи (в тому числі й власної суспільної природи). Працювати, обробляти природний матеріал, створювати речі, яких немає в природі – ось що принципово відрізняє життя людини від життєдіяльності тварини і визначає всі інші відмінності. Але досі сама праця в суспільному масштабі була невільною. Усі докомуністичні способи виробництва відрізнялися своєю специфічною формою невільної праці. Комунізм же передбачає повне подолання цієї неволі. Тому його основою може бути лише вільна праця «без норми та винагороди», як «найперша життєва потреба» людини.

Найпростіша, безпосередньо колективна форма здійснення такої праці (а краще сказати, такої діяльності, оскільки слово «праця» певною мірою зберігає негативне значення) є тією самою «клітиною» комунізму, в якій закладено подальший розвиток комуністичного суспільства у всіх його необхідних визначеннях.

Вільна діяльність та її безпосередньо-колективна форма є реальністю в наш час, так само як реальністю ще до появи капіталізму були товарне виробництво та товар. І перехід до нового суспільного ладу відбудеться, коли вільна діяльність людини стане універсальною формою, так само, як колись товаром. З розмаїття товарів на ринку, що мають різні властивості, неможливо безпосередньо зрозуміти, що таке товар і товарне виробництво. З різноманітності видів і способів організації комуністичної діяльності неможливо зрозуміти її суть і те, що таке її безпосередньо-колективна форма. Тим більше, що слова можуть затьмарювати суть справи. Безпосередньо-колективна форма вільної діяльності може існувати навіть тоді, коли людина здійснює її зовсім одна. Необхідно виділити ті визначення, які є необхідними для неї.

У безпосередньо-колективній формі вільна діяльність здійснюється за критеріями свободи, істини, добра та краси, а отже, і найвищих людських інтересів, тобто це форма самоцінної чуттєвої діяльності. Вільна діяльність обов’язково передбачає єдність усіх трьох компонентів: чуттєвості, моральності та мислення, і їх піднесення до рівня свободи. У процесі такої діяльності здійснюється утвердження людини в усій її повноті незалежно від того, що саме робиться і що саме (істина, добро чи краса) виходить на передній план.

Така безпосередньо-колективна форма – це лише зовнішній вигляд самої діяльності разом із відповідним їй змістом. У повсякденній мові таку діяльність часто називають соціальною творчістю або просто творчістю. Форма є важливою, оскільки розвиток відбувається та виражається у необхідному русі форми. Зміст завжди породжує відповідну форму та здійснюється через неї. Тому протиріччя змісту спочатку сприймаються свідомістю як протирічливість форми. У її розгортанні та через її розгортання (зміну форм) насправді відбувається рух змісту. З боку форми в діяльності важливими є такі моменти, як характер зв’язку між людьми, мета, суспільне буття предметів і засобів праці.

Безпосередня колективність у вільній діяльності проявляється, перш за все, в орієнтації на інших людей. У безпосередньо колективних формах діяльності людина завжди виступає як мета та найвища цінність. Усе, що робиться, робиться безпосередньо для розвитку інших людей. Тому продукт такої праці ніколи не набуває самостійно-речової форми (форми товару). Продукт такої праці не може протистояти живим людям як чужа їм річ. Іншими словами, ця діяльність спрямована на розвиток суспільства та покращення життя інших людей, здійснюється для них, а отже, і для себе, як представника цього суспільства. Те, що зараз прийнято називати матеріальними цінностями у вільній комуністичній діяльності, може слугувати лише засобом.

Що стосується характеру зв’язку між людьми, то його загалом можна визначити як товариство, або, що те саме, вільна асоціація.

Цей характер взаємозв’язку також вказує на засоби такої діяльності, оскільки у вільній асоціації жоден з учасників не може бути засобом для здійснення чужої волі та чужих йому інтересів. Навіть якщо хтось, завдяки своїм здібностям і навичкам, є провідним у колективі, інші не сліпо виконують чужу волю, а співпрацюють у здійсненні спільної справи, яка є однаково важливою та необхідною для всіх її учасників.

Характер зв’язку також вказує на суспільну форму існування засобів і предмета комуністичної діяльності. Таким чином, засобами можуть виступати лише речі та процеси, відмінні від людини. Суспільна форма існування засобів і предметів праці передбачає їх безпосередньо суспільне використання. Саме в цьому сенсі Маркс говорив про перехід від управління людьми до управління речами.

З вищесказаного випливає, що такі форми людської діяльності позбавлені будь-якого зовнішнього примусу, як прямого, так і непрямого (винагороди за працю), оскільки ні сам процес, ні результат комуністичної діяльності не є байдужими для того, хто її здійснює. Людина щоразу зацікавлена у своїй праці безпосередньо, на відміну від ситуації, яка є звичайною для сучасного суспільства, коли сама діяльність виступає лише засобом для отримання грошей.

Звісно, безпосередня колективність полягає ще й у тому, що праця кожного члена суспільства спирається на працю інших людей, у тому, що предмети, засоби та способи праці є продуктом, створеним іншими людьми, а його власна праця, у свою чергу, є невід’ємною частиною всієї суспільної праці та надає матеріал для майбутньої діяльності. Але це характерно для будь-якої праці, що здійснюється в будь-якій суспільній формі, а не лише для комуністичної діяльності, тому це не стосується специфіки комунізму.

Під безпосередньо-колективною формою вільної діяльності ми розуміємо об’єднання людей у процесі діяльності, де метою кожного є розвиток усіх, а засобами для її досягнення виступають речі та природні процеси, щодо яких відбувається вільне та рівне вираження волі.

Безпосередньо-колективна форма діяльності – це не просто зовнішня суспільна форма, яка може існувати, а може й ні, не впливаючи при цьому на зміст діяльності, а навпаки, – це її змістовна форма. Без неї взагалі неможливе здійснення вільної діяльності, що відповідає критеріям істини, добра та краси.\footnote{Йдеться не про поширене уявлення, де краса пов’язується зі смаком, добро – з думками про те, що є добре і що є погано, і ці уявлення можуть бути різними, а лише істині приписується якесь об’єктивне значення. Категорії «істини», «добра» та «краси» розроблялися протягом тисячоліть як об’єктивні ідеальні характеристики людської діяльності у світі, що розгортаються (розвиваються) в цій діяльності, а отже, і як об’єктивні характеристики світу (матерії). Іншими словами, якщо коротко, коли йдеться про діяльність, яка оцінюється за критеріями істини, добра та краси, то маються на увазі «точки розвитку» людської чуттєвості, пізнання та моралі, розвиток відносин між людьми та відносин людини з природою, розглянутих з ідеальної точки зору.} Розвиток і метаморфози такої форми передбачають збереження вищезазначених моментів, подібно до того, як у метаморфозах товарної форми продуктів праці завжди зберігається вартість (те, що робить товар саме товаром, а не просто продуктом праці).

Що ж до змісту такої діяльності, то вона завжди є революційною. У цьому полягає її внутрішнє протиріччя. У вільній комуністичній діяльності постійно закладається як її внутрішня межа, так і вихід за ці межі. За допомогою історично досягнутих, а отже, історично обмежених рубежів діяльності, змінюється сама ця діяльність, закладається її новий зміст (завдяки розширенню горизонтів освоєння матерії) та нова форма. Це відбувається не незалежно від учасників процесу праці, обхідним шляхом, як це відбувається в процесі матеріального виробництва. Навпаки, ця зміна відбувається цілеспрямовано. Тому комунізм, як оволодіння людиною власною діяльністю, виявляє внутрішню суперечливість людської діяльності взагалі, а отже, і внутрішню суперечливість матеріального світу, в якому людина здійснює свою діяльність. Тому для науки про комунізм найважливіше методологічне значення має зауваження Маркса, висловлене в першій тезі про Фейєрбаха: «Головний недолік усього попереднього матеріалізму – включно з фейєрбахівським – полягає в тому, що предмет, реальність, чуттєвість розглядаються лише у формі об’єкта, або у формі споглядання, а не як людська чуттєва діяльність, практика, не суб’єктивно. Звідси й сталося, що діяльна сторона, на противагу матеріалізму, розвивалася ідеалізмом, але лише абстрактно, оскільки ідеалізм, звісно, не знає дійсньої, чуттєвої діяльності як такої. Фейєрбах хоче мати справу з чуттєвими об’єктами, які справді відрізняються від уявних понять про те, що є добре і що є погано, які можуть бути різними, і лише істині надається якесь об’єктивне значення. Категорії «істини», «добра» і «краси» розроблялися протягом тисячоліть як об’єктивні ідеальні характеристики людської діяльності у світі, що розгортаються (розвиваються) в цій діяльності, а тому й об’єктивні характеристики світу (матерії). Іншими словами, якщо коротко, коли йдеться про діяльність за мірками істини, добра і краси, то маються на увазі «точки розвитку» людської чуттєвості, пізнання і моральності, про розвиток ставлення людини до людини і людини до природи, взятих з боку ідеального виміру об’єктів, але найголовніше, він розглядає людську діяльність не як предметну діяльність. У цій заяві підкреслюється не лише роль ідеалізму у формуванні теорії діяльності. Хоча ця роль справді величезна і важлива, настільки, що згодом Ленін говорив про важливість розвитку «матеріалістичних друзів гегелівської діалектики». Матеріалістично зрозумілі категорії діалектики як «мільйони разів повторювані форми людської діяльності» (Ленін), їх внутрішня суперечливість і рух є необхідним інструментом розвитку теорії комунізму саме тому, що в них відображені суттєві моменти людської діяльності взагалі, у них відображено те, без чого людська діяльність в принципі не може здійснюватися. Але в марксистському тезі також робиться акцент на матеріальності людської діяльності.

Підкреслимо ще раз, що вільна діяльність – це не «відпочинок» у сучасному розумінні цього слова, не безцільне проведення часу, а діяльність, що відповідає критеріям свободи. Колись філософія, в особі Бенедикта Спінози, визначила свободу як усвідомлену необхідність. Справжня свобода визначається діями, що відповідають цій необхідності.

У вільний час, відповідно до історичної необхідності, люди розвивають культуру в найширшому сенсі цього слова. Така діяльність не є вигадкою комуністів у фантастичних мріях. Вона існувала в минулому і існує зараз. Але вона існувала і існує в дуже жорстких рамках і є надбанням лише небагатьох людей. Абсолютна більшість членів суспільства не знає вільної діяльності. Це відбувалося і відбувається, зокрема, тому, що потрібно не просто мати час, а й бути розвиненим настільки, щоб бути в змозі здійснювати таку діяльність. Умови для такого розвитку були і залишаються в сучасному товарному суспільстві досить обмеженими.

Отже, щоб рухатися в напрямку комунізму, потрібно виробляти комунізм, тобто вільну діяльність, точніше, вільний час і відповідний рівень розвитку членів суспільства, які могли б її здійснювати.

Отже, і розподіл тут, насамперед, – це розподіл вільного часу та можливостей для засвоєння людської культури. Розподіл продуктів праці необхідний лише настільки, наскільки він забезпечує вільний час. У цьому ключі слід ставити питання про ефективність суспільного виробництва в перехідний період до комунізму. Вільний час при капіталізмі є виробничими витратами, його не враховують, не цінують, і тому він часто так і не стає справді вільним, не використовується за призначенням. Призначення вільного часу – бути використаним для постійного розвитку суспільних відносин (або розвитку діяльності, що, по суті, одне й те саме).

Отже, для комунізму надзвичайно важливо, щоб звичка, яка стала надбанням усіх, ставитися до власного вільного часу та діяльності, що здійснюється в цей час, як до суспільного багатства, яке не підлягає приватному розбазарюванню, була сформована. Загалом, сила звички та колективного формування звички має тут величезне значення, хоча б тому, що бути вільним, тобто діяти за принципами свободи, завжди дуже важко. Це означає перебувати на передовій суперечностей, через які здійснюється розвиток суспільства. Це не лише важка праця, але й трагедія в тому сенсі, що потрібно постійно долати себе, заперечувати себе минулого. А колективна суспільна свобода – це колективна суспільна трагедія та її постійне розв’язання. Комуністичне суспільство – це не суспільство загальної лінощів і бездіяльності, а суспільство загальної вільної діяльності, яка є для людей першою життєвою потребою.

Комунізм реалізується через необхідну зміну безпосередньо колективних форм людської діяльності, спрямованої на зміну самої цієї діяльності (розвиток самого життя людей). Досягнуті форми, прагнучи вийти за свої межі, створюють нові. Тому цю зміну не можна уявляти як постійно повторювані, що йдуть один за одним, однакові метаморфози (як ми бачимо це в русі капіталу). Внутрішня логіка руху безпосередньо колективних форм вільної діяльності, необхідність їхньої зміни визначається не просто історично досягнутим способом виробництва та відповідними йому виробничими відносинами, а свободою як необхідністю самої матерії, що розкривається в міру освоєння матерії.

У найзагальнішому вигляді закон комунізму можна сформулювати як закон свободи або закон самодетермінації людської діяльності: вільна діяльність членів суспільства, що здійснюється як усвідомлений розвиток змісту та форм цієї діяльності, визначає розвиток усього суспільства. Або, іншими словами, свобода, реалізуючись, прагне до самозростання.

Звісно, розуміння дії цього закону потребує постійного уточнення з використанням усіх наявних знань як про природу, так і про суспільство. Тому теорія комунізму може за відповідних умов відігравати об’єднуючу роль для всіх наук, вказуючи на реальний зв’язок між ними, як зв’язок самої людської практики.

Існуючі сьогодні безпосередньо-колективні форми вільної людської діяльності, як «зачатки комунізму», є не просто свідченням його можливості, а й його реальним існуванням. Важливо усвідомити, що комуністичний спосіб виробництва вже існує, йдеться лише про те, щоб він став домінуючим і поширився на все суспільство.

У своєму русі безпосередньо-колективні форми вільної людської діяльності, як «зачатки комунізму», можуть розгорнутися в комунізм як суспільний уклад, тобто як спосіб існування суспільства. У цьому сенсі одним із найцінніших положень у теоретичній спадщині Леніна є закликати виявляти та плекати «зачатки» комунізму всюди, де вони є.
\end{document}
